\documentclass[12pt]{article}
\usepackage[utf8]{inputenc}
\usepackage[T1]{fontenc}
\usepackage{hyperref}
\title{Puerto Rican Sign Language Language Resources for K-12 Curriculum Development}
\author{Compiled by Claude Sonnet 5.5 (claude-sonnet-5-5) for Andr\'as Kornai}
\date{}
\begin{document}
\maketitle

\section{Introduction}
Puerto Rican Sign Language (LSP; ISO 639-3: psl; Glottocode: puer1237) is the sign language of the Deaf community of Puerto Rico. The search results were mostly Puerto Rican legislative documents and news items in Spanish, plus one Department of Education resource page. No dictionary, grammar or signed-language teaching material was found.

\section{Primary Educational Resources}
\begin{itemize}
\item Puerto Rico Department of Education digital resource page (per search summary, linked to a 2024 guide on providing and integrating LSP in the curriculum; the page itself was not opened): \url{https://dedigital.dde.pr/mod/resource/view.php?id=3644}
\end{itemize}

\section{Linguistic and Reference Materials}
No grammar, dictionary or linguistic study was found in the search results.

\section{Religious and Community Resources}
\begin{itemize}
\item Opinion piece on sign language assistance in Puerto Rico, El Nuevo D\'ia: \url{https://www.elnuevodia.com/opinion/punto-de-vista/ayuda-o-monopolio-de-la-lengua-de-senas-en-puerto-rico}
\item News on Deaf residents during the 2017 emergency, Primera Hora: \url{https://www.primerahora.com/noticias/puertorico/nota/sordossesientenolvidadosenmediodelaemergenciaesunafaltaderespeto-1386480/}
\item News on the number of Deaf people in Puerto Rico, Noticel: \url{https://www.noticel.com/ahora/cuidate/20160505/puerto-rico-desconoce-cuantos-sordos-viven-en-el-pais}
\end{itemize}

\section{Digital and Audio Resources}
No video or digital resources in LSP were found in the search results.

\section{Curriculum Development Context}
Policy documents from the Puerto Rico Senate bear on recognition and schooling: a 2025 bill (P.~del S.~618) to recognize LSP as the natural language of the Deaf community and its veto message, a 2021 bill to add it to the official languages, and 2017 bills on sign language courses in schools. Whether any became law was not confirmed. The bill texts state that American Sign Language interpreting in emergencies caused confusion because LSP differs from ASL, which supports a curriculum based on LSP itself taught by Deaf signers.
\begin{itemize}
\item P.~del S.~618 (2025): \url{https://senado.pr.gov/document_vault/legislative_measures/5662/document/ps0618-25.pdf}
\item Veto message on P.~del S.~618: \url{https://senado.pr.gov/document_vault/vetos/ps0618e-25.pdf}
\item P.~del S.~481 (2021): \url{https://senado.pr.gov/document_vault/legislative_measures/111/document/ps0481-21.pdf}
\item P.~del S.~445 (2017): \url{https://senado.pr.gov/document_vault/legislative_measures/1784/document/ps0445-17.pdf}
\item P.~del S.~606 (2017): \url{https://senado.pr.gov/document_vault/legislative_measures/1619/document/ps0606-17.pdf}
\end{itemize}

\end{document}